% !TeX root = ../main.tex
\section{Исследование нелинейных свойств усилителя и измерение точки децибельной компрессии ($P_{1\text{dB}}$)}

\subsection{Теоретические сведения и методика проведения измерений}

При увеличении мощности входного сигнала амплитудная характеристика (АХ) реального усилителя перестает быть линейной вследствие физических ограничений активных полупроводниковых структур (насыщения по току коллектора/стока и ограничения динамического диапазона напряжением питания). 

Параметром, определяющим верхнюю границу линейного динамического диапазона усилителя, является \textbf{точка децибельной компрессии} ($P_{1\text{dB}}$). Она соответствует уровню мощности, при котором реальный коэффициент усиления каскада падает на 1~дБ по сравнению с коэффициентом усиления на начальном малосигнальном (линейном) участке характеристики:
\begin{equation}
OP_{1\text{dB}} = IP_{1\text{dB}} + K_1 - 1\text{ дБ},
\end{equation}
где $IP_{1\text{dB}}$ --- входная мощность компрессии в 1~дБ, $OP_{1\text{dB}}$ --- выходная мощность компрессии в 1~дБ, а $K_1$ --- малосигнальный коэффициент усиления.

\textbf{Порядок выполнения экспериментальных исследований:}
\begin{enumerate}
    \item \textbf{Калибровка потерь выходного тракта ($IL_2$):}
    \begin{itemize}
        \item собрана измерительная схема без исследуемого усилителя (рис.~24 методических указаний): Генератор ($+10\text{ дБм}$) $\to$ аттенюатор 20~дБ $\to$ Кабель 1 $\to$ адаптер $\text{SMA}(f)\text{--}\text{SMA}(f)$ $\to$ аттенюатор 10~дБ $\to$ Кабель 2 $\to$ Анализатор спектра; зафиксированы калибровочные уровни $P_1(f)$;
        \item из схемы исключены выходной аттенюатор 10~дБ и Кабель 2 (рис.~25 методических указаний); зафиксированы уровни $P_2(f)$;
        \item суммарные потери выходного измерительного плеча $IL_2 = P_2 - P_1$ составили порядка $11.0\text{ дБ}$.
    \end{itemize}
    \item \textbf{Снятие амплитудных характеристик (рис.~26 методических указаний):} в измерительный тракт включен исследуемый усилитель с подачей напряжения питания. Мощность генератора сигналов варьировалась в диапазоне $P_{\text{вх}} \in [-10; +18]\text{ дБм}$ с шагом 2~дБ на частотах $0.5\text{ ГГц}$, $1.0\text{ ГГц}$ и $2.5\text{ ГГц}$. На анализаторе спектра регистрировался уровень мощности $P_{\text{изм}}$.
    \item \textbf{Восстановление реальной мощности усилителя:} истинная выходная мощность $P_{\text{вых}}$ определялась компенсацией потерь выходного плеча:
    \begin{equation}
    P_{\text{вых}} = P_{\text{изм}} + IL_2.
    \end{equation}
    \item \textbf{Определение точки компрессии:} строилась линейная асимптота по начальным точкам малого сигнала ($P_{\text{вх}} \in [-10; -6]\text{ дБм}$). Пороговое значение децибельного сжатия определялось из условия отклонения экспериментальной кривой от линейной аппроксимации ровно на 1~дБ ($\Delta = P_{\text{вых, идеал}} - P_{\text{вых}} = 1.0\text{ дБ}$).
\end{enumerate}

\subsection{Результаты экспериментальных исследований}

Экспериментальные данные измерений выходной мощности на анализаторе спектра и рассчитанные значения приведенной мощности $P_{\text{вых}}$ для частот $0.5\text{ ГГц}$, $1.0\text{ ГГц}$ и $2.5\text{ ГГц}$ приведены в таблицах~\ref{tab:p1db_1ghz} и~\ref{tab:p1db_multi}.

\begin{table}[H]
\centering
\caption{Экспериментальные данные и результаты расчета амплитудной характеристики ($f = 1\text{ ГГц}$, $IL_2 = 11.02\text{ дБ}$)}
\label{tab:p1db_1ghz}
\small
\begin{tabular}{|c|c|c|c|c|}
\hline
$P_{\text{вх}}$, дБм & $P_{\text{изм}}$, дБм & $P_{\text{вых}}$, дБм & $P_{\text{вых, идеал}}$, дБм & Сжатие $\Delta$, дБ \\
\hline
-10 & -18.39 & -7.37 & -7.33 & 0.04 \\
-8  & -16.49 & -5.47 & -5.55 & -0.08 \\
-6  & -14.83 & -3.81 & -3.77 & 0.04 \\
-4  & -13.20 & -2.18 & -1.99 & 0.19 \\
-2  & -12.02 & -1.00 & -0.21 & 0.79 \\
0   & -10.93 &  0.09 &  1.57 & 1.48 \\
2   & -10.05 &  0.97 &  3.35 & 2.38 \\
4   & -9.34  &  1.68 &  5.13 & 3.45 \\
6   & -8.64  &  2.38 &  6.91 & 4.53 \\
8   & -8.10  &  2.92 &  8.69 & 5.77 \\
10  & -7.66  &  3.36 & 10.47 & 7.11 \\
12  & -7.23  &  3.79 & 12.25 & 8.46 \\
14  & -6.86  &  4.16 & 14.03 & 9.87 \\
16  & -6.41  &  4.61 & 15.81 & 11.20 \\
18  & -6.07  &  4.95 & 17.59 & 12.64 \\
\hline
\end{tabular}
\end{table}

\begin{table}[H]
\centering
\caption{Экспериментальные данные и результаты расчета АХ для частот $0.5\text{ ГГц}$ и $2.5\text{ ГГц}$ ($IL_2 = 11.0\text{ дБ}$)}
\label{tab:p1db_multi}
\small
\begin{tabular}{|c|c|c|c|c|c|c|}
\hline
\multirow{2}{*}{$P_{\text{вх}}$, дБм} & \multicolumn{3}{c|}{$f = 0.5\text{ ГГц}$} & \multicolumn{3}{c|}{$f = 2.5\text{ ГГц}$} \\
\cline{2-7}
 & $P_{\text{изм}}$, дБм & $P_{\text{вых}}$, дБм & $\Delta$, дБ & $P_{\text{изм}}$, дБм & $P_{\text{вых}}$, дБм & $\Delta$, дБ \\
\hline
-10 & -16.82 & -5.82 & 0.02  & -22.75 & -11.75 & -0.02 \\
-8  & -14.98 & -3.98 & -0.04 & -20.98 & -9.98  &  0.04 \\
-6  & -13.27 & -2.27 & 0.02  & -19.09 & -8.09  & -0.02 \\
-4  & -11.80 & -0.80 & 0.33  & -17.40 & -6.40  &  0.12 \\
-2  & -10.64 &  0.36 & 0.94  & -15.65 & -4.65  &  0.20 \\
0   &  -9.67 &  1.33 & 1.75  & -14.48 & -3.48  &  0.86 \\
2   &  -8.77 &  2.23 & 2.62  & -13.43 & -2.43  &  1.64 \\
4   &  -7.86 &  3.14 & 3.49  & -12.34 & -1.34  &  2.38 \\
6   &  -6.99 &  4.01 & 4.39  & -11.33 & -0.33  &  3.20 \\
8   &  -6.15 &  4.85 & 5.33  & -10.42 &  0.58  &  4.12 \\
10  &  -5.44 &  5.56 & 6.39  &  -9.33 &  1.67  &  4.86 \\
12  &  -4.83 &  6.17 & 7.56  &  -8.44 &  2.56  &  5.80 \\
14  &  -4.63 &  6.37 & 9.13  &  -7.78 &  3.22  &  6.97 \\
16  &  -4.46 &  6.54 & 10.74 &  -7.29 &  3.71  &  8.31 \\
18  &  -4.41 &  6.59 & 12.46 &  -6.99 &  4.01  &  9.84 \\
\hline
\end{tabular}
\end{table}

На рисунках~\ref{fig:p1db_1ghz}--\ref{fig:p1db_multi} приведены построенные характеристики сжатия динамического диапазона усилителя с нанесенными идеальными линейными асимптотами и расчетными точками компрессии $P_{1\text{dB}}$.

\begin{figure}[H]
    \centering
    \includegraphics[width=0.88\textwidth]{Р1дб.png}
    \caption{Характеристика сжатия динамического диапазона и точка децибельной компрессии ($f = 1.0\text{ ГГц}$)}
    \label{fig:p1db_1ghz}
\end{figure}

\begin{figure}[H]
    \centering
    \begin{minipage}[b]{0.48\textwidth}
        \centering
        \includegraphics[width=\textwidth]{Р1дб0,5.png}
        \caption{Характеристика компрессии при $f = 0.5\text{ ГГц}$}
        \label{fig:p1db_05ghz}
    \end{minipage}
    \hfill
    \begin{minipage}[b]{0.48\textwidth}
        \centering
        \includegraphics[width=\textwidth]{Р1дб2,5.png}
        \caption{Характеристика компрессии при $f = 2.5\text{ ГГц}$}
        \label{fig:p1db_25ghz}
    \end{minipage}
    \label{fig:p1db_multi}
\end{figure}

Сводные результаты определения координат точки компрессии по результатам сплайн-интерполяции приведены в таблице~\ref{tab:p1db_summary}.

\begin{table}[H]
\centering
\caption{Сводные параметры точки компрессии $P_{1\text{dB}}$ на исследуемых частотах}
\label{tab:p1db_summary}
\begin{tabular}{|c|c|c|c|}
\hline
\begin{tabular}[c]{@{}c@{}}Частота $f$,\\ ГГц\end{tabular} & 
\begin{tabular}[c]{@{}c@{}}Входная точка\\ $IP_{1\text{dB}}$, дБм\end{tabular} & 
\begin{tabular}[c]{@{}c@{}}Выходная точка\\ $OP_{1\text{dB}}$, дБм\end{tabular} & 
\begin{tabular}[c]{@{}c@{}}Мощность насыщения\\ $P_{\text{sat}}$, дБм\end{tabular} \\
\hline
0.5 & -1.82 & +0.46 & $\approx +6.6$ \\
1.0 & -1.36 & -0.64 & $\approx +5.0$ \\
2.5 & +0.37 & -3.29 & $\approx +4.0$ \\
\hline
\end{tabular}
\end{table}

\subsection{Анализ результатов и выводы}

\begin{enumerate}
    \item \textbf{Зависимость точки компрессии по выходу ($OP_{1\text{dB}}$) от частоты:} С увеличением рабочей частоты от 0.5~ГГц до 2.5~ГГц выходная мощность децибельного сжатия $OP_{1\text{dB}}$ монотонно снижается (с $+0.46\text{ дБм}$ на частоте 0.5~ГГц до $-0.64\text{ дБм}$ на 1.0~ГГц и $-3.29\text{ дБм}$ на 2.5~ГГц). Соответственно уменьшается и предельная мощность насыщения каскада ($P_{\text{sat}}$ падает с $+6.6\text{ дБм}$ до $+4.0\text{ дБм}$). Спад обусловлен деградацией усилительных свойств транзистора, ростом паразитных высокочастотных потерь в полупроводниковой структуре и цепях межкаскадного согласования.
    \item \textbf{Поведение точки компрессии по входу ($IP_{1\text{dB}}$):} На частоте 2.5~ГГц точка компрессии по входу смещается в область больших мощностей ($IP_{1\text{dB}} = +0.37\text{ дБм}$ против $-1.82\text{ дБм}$ на 0.5~ГГц). Это объясняется тем, что из-за меньшего коэффициента усиления $S_{21}$ на высоких частотах для раскачки выходного каскада до ограничения напряжения и тока требуется подача существенно большей входной мощности от генератора.
    \item \textbf{Динамический диапазон каскада:} Экспериментально установленные значения $P_{1\text{dB}}$ определяют верхний предел мощности входного сигнала. Для сохранения линейности радиотракта и предотвращения интермодуляционных искажений номинальная рабочая мощность должна выбираться с запасом не менее 10--12~дБ относительно точки компрессии.
\end{enumerate}